% Propietario conservado: /Users/ruben/Documents/ChatGPT/jueces y controles/REVISION_ARGUMENTAL_APERTURAS_CIERRES_20260915/delivery/ARTICLE_V_ES_ARGUMENTAL_20260915/manuscrito/sections/v_entropia_sectorial.tex
% RESULTADO_RECUPERADO desde II REV09. Véase MAPA_PROCEDENCIA.json.
% Selección de líneas y cambios mecánicos explícitos; ninguna prueba nueva.
\subsection{Equation of state for sectorial information}
\label{v:v:ant:sec:ii-ecuacion-informacion}

The incidence already constructed provides eighteen tritic factors in
each sector. This multiplicity allows the relation between area,
the occupation distribution, and the entropy of the preceding
purification to be written explicitly. We retain the reduced state defined there
and abbreviate $\Delta=\Delta_{\rm mod}$. For $m\in\{90,120\}$, put
\begin{equation}
 x_m=e^{-m\Delta},\qquad z_m=1+x_m+x_m^2,\qquad
 f_m=\frac{x_m+2x_m^2}{z_m},\qquad
 v_m=\frac{x_m+4x_m^2}{z_m}-f_m^2.
 \label{v:v:ant:eq:ii-sector-momentos}
\end{equation}
Here $f_m$ and $v_m$ are, respectively, the mean and variance of
the occupation of one factor in sector $m$. The parameter $\Delta$ preserves
the status and normalization of the previously specified reduced
state.

\subsubsection{Bipartite quantum purification of the incidence sectors}
\label{v:v:ant:sec:ii-area-purificacion-producto}

The preceding Schmidt decomposition is made concrete on each link.
For a finite real weight \(w\), define
\[
 z(w)=1+e^{-w}+e^{-2w},\qquad
 \rho(w)=\frac1{z(w)}\sum_{n=0}^2e^{-wn}|n\rangle\langle n|.
\]
The sectorial weights are \(w_m=m\Delta\), with \(m=90,120\);
hence \(z(w_m)=z_m\).

\begin{proposition}[Thermofield-double purification of sectorial incidence]
\label{v:v:ant:prop:ii-area-purificacion-producto}
The normalized vector
\begin{equation}
 |\operatorname{TFD}(w)\rangle
 =\frac1{\sqrt{z(w)}}\sum_{n=0}^2
       e^{-wn/2}|n\rangle_A\otimes|n\rangle_{\bar A}
 \label{v:v:ant:eq:ii-area-tfd-enlace}
\end{equation}
purifies \(\rho(w)\). The product
\begin{equation}
 |\Psi_\Delta\rangle
 =|\operatorname{TFD}(90\Delta)\rangle^{\otimes18}
  \otimes|\operatorname{TFD}(120\Delta)\rangle^{\otimes18}
 \label{v:v:ant:eq:ii-area-tfd-producto}
\end{equation}
has reduction \(\rho_A(\Delta)\) on the thirty-six original
links. Its entanglement entropy for the bipartition
\(A\mid\bar A\) is \(s(\Delta)\).
\end{proposition}
\begin{proof}
The squared norm is \(z(w)^{-1}\sum_{n=0}^2e^{-wn}=1\).
In the partial trace, orthogonality of the second factor eliminates
terms with distinct indices; the remaining terms form \(\rho(w)\).
The partial trace commutes with the tensor product. One thus obtains
\(\rho(90\Delta)^{\otimes18}\otimes\rho(120\Delta)^{\otimes18}\),
which coincides with the normalized state defined previously.
Equality between the entropy of a reduction and the entanglement of
this purification is read directly from its Schmidt coefficients.
\end{proof}

The two factors of the bipartition are determined within the
construction. The physical realization of the second factor preserves that
subsystem identification and its algebra of observables.

\subsubsection{Occupation moments and information content}

\begin{proposition}[Mean area and entropy of the sectorial purification]
In this finite family of states,
\begin{align}
 \langle N\rangle_\Delta&=18(f_{90}+f_{120}),\\
 \langle D\rangle_\Delta&=18(90f_{90}+120f_{120}),\\
 \frac{\langle\widehat{\mathcal A}\rangle_\Delta}{\ell_P^2}
 &=18\gamma_{\HMT}a_0(f_{90}+f_{120}),
 \label{v:v:ant:eq:ii-area-media}\\
 s(\Delta)&=18\log z_{90}+18\log z_{120}
                  +\Delta\langle D\rangle_\Delta.
 \label{v:v:ant:eq:ii-entropia-sectorial}
\end{align}
The dimensional entropy is $S(\Delta)=k_Bs(\Delta)$, with the
thermal transducer of Section~\ref{v:v:ant:sec:ii-kb-aut-desarrollo}.
\end{proposition}
\begin{proof}
In each factor, the normalized weights are
$(1,x_m,x_m^2)/z_m$. Their first moment gives $f_m$, and their second
moment gives the first fraction in $v_m$. Tensor factorization
and additivity of the number operators produce the first three
identities. Finally,
$\log\rho_A=-(\log Z)I-\Delta D$, with
$Z=z_{90}^{18}z_{120}^{18}$. Taking the trace against $-\rho_A$
gives~\eqref{v:v:ant:eq:ii-entropia-sectorial}. The purification specified
previously has this same entanglement entropy.
\end{proof}

Loss of uniformity has its own quantitative reading.
In dimension \(d_\partial=3^{36}\), let
\(\tau_\partial=I/d_\partial\). Define
\begin{equation}
 I_{\rm or}(\Delta)=36\log3-s(\Delta)
 =D\bigl(\rho_A(\Delta)\Vert\tau_\partial\bigr)\ge0.
 \label{v:v:ant:eq:ii-area-deficit-informacional}
\end{equation}
The equality is obtained by substituting
\(\log\tau_\partial=-36\log3\,I\). If \(p_\nu\) are the
eigenvalues of \(\rho_A\), the convexity of \(t\log t\) gives
\(\sum_\nu p_\nu\log(d_\partial p_\nu)\ge0\), with equality
exactly for the uniform state. For this family, this occurs at
\(\Delta=0\). The deficit is dimensionless information; the areal
quantum \(\gamma_{\HMT}a_0\) is another reading. Each expression retains
its definition even though their evaluations may be close.

\begin{proposition}[Information response and boundary fluctuations]
For every finite real $\Delta$,
\begin{align}
 \frac{\dd\langle D\rangle_\Delta}{\dd\Delta}
 &=-\operatorname{Var}_\Delta(D)
   =-18(90^2v_{90}+120^2v_{120}),\\
 \frac{\dd s}{\dd\Delta}
 &=-\Delta\operatorname{Var}_\Delta(D),
 \label{v:v:ant:eq:ii-respuesta-entropica}\\
 \frac{\dd\langle\widehat{\mathcal A}\rangle_\Delta}{\dd\Delta}
 &=-18\gamma_{\HMT}a_0\ell_P^2(90v_{90}+120v_{120}).
 \label{v:v:ant:eq:ii-respuesta-areal}
\end{align}
\end{proposition}
\begin{proof}
The sum defining $Z$ is finite. Differentiating it gives
$\partial_\Delta\log Z=-\langle D\rangle_\Delta$ and
$\partial_\Delta^2\log Z=\operatorname{Var}_\Delta(D)$.
The sectors and their factors are independent in the product state;
their variances therefore add. On differentiating
\eqref{v:v:ant:eq:ii-entropia-sectorial}, the terms
$-\langle D\rangle_\Delta$ and $\langle D\rangle_\Delta$
cancel. For area, use
$\partial_\Delta f_m=-mv_m$ in~\eqref{v:v:ant:eq:ii-area-media}.
\end{proof}

These equations relate three levels of the same construction.
Incidence fixes multiplicities and sectoral weights; the
reduced state fixes their probabilities; fluctuations of those
occupations determine the response of area and entropy.
The scale $\gamma_{\HMT}a_0\ell_P^2$ converts occupation into area,
while $k_B$ converts information into dimensional entropy.
In section~\ref{v:v:ant:sec:gravity}, $\ell_P^2=G\hbar/c^3$ makes explicit
the dependence of that areal unit on action and gravitation.
The family of reduced states and its bipartition remain specified
in this development; any horizon realization must also transport
those two data.

\subsection{Occupation inversion and entropic invariants}
\label{v:v:ant:sec:ii-inversion-sectorial}

Let $\mathcal R$ be the unitary involution that exchanges $|0\rangle$
and $|2\rangle$ and fixes $|1\rangle$ in each of the thirty-six
tritic factors. Its action on occupations preserves the entire
microscopic configuration through a bijective change of its labels.

\begin{theorem}[Preserved entropy and complementary area]
The preceding involution satisfies
\begin{align}
 \mathcal R N\mathcal R^*&=72I-N,&
 \mathcal R D\mathcal R^*&=7560I-D,\\
 \rho_A(-\Delta)&=\mathcal R\rho_A(\Delta)\mathcal R^*,&
 s(-\Delta)&=s(\Delta),
 \label{v:v:ant:eq:ii-inversion-entropia}\\
 \langle\widehat{\mathcal A}\rangle_{-\Delta}
 +\langle\widehat{\mathcal A}\rangle_\Delta
 &=72\gamma_{\HMT}a_0\ell_P^2.&&
 \label{v:v:ant:eq:ii-area-complementaria}
\end{align}
The relative entropy between the two orientations is
\begin{equation}
 D\bigl(\rho_A(\Delta)\Vert\rho_A(-\Delta)\bigr)
 =\Delta\bigl(7560-2\langle D\rangle_\Delta\bigr).
 \label{v:v:ant:eq:ii-distancia-orientaciones}
\end{equation}
\end{theorem}
\begin{proof}
Each local occupation $n$ transforms into $2-n$. Thus,
$N_m\mapsto36I-N_m$ and
$D\mapsto36(90+120)I-D=7560I-D$.
Furthermore, $z_m(-\Delta)=e^{2m\Delta}z_m(\Delta)$,
whence $Z(-\Delta)=e^{7560\Delta}Z(\Delta)$.
Substituting both identities into the normalized state proves
unitary conjugation. Spectral invariance proves equality
of entropies; transformation of $N$ proves the area sum.
Finally, subtracting the two state logarithms and taking
their expectation in $\rho_A(\Delta)$ yields
\eqref{v:v:ant:eq:ii-distancia-orientaciones}.
\end{proof}

At $\Delta=0$ the reduced state is uniform and
$s(0)=36\log3$. In the limit $\Delta\to+\infty$, the zero-occupation
configuration survives; the entropy tends to zero.
Inversion takes that limit to the maximum-occupation configuration.
At both extremes the entropy is the same and the area distinguishes the
configurations. For $\Delta>0$, the response
\eqref{v:v:ant:eq:ii-respuesta-entropica} is strictly negative because
the state has full rank and $D$ has more than one eigenvalue.
The law thus preserves the difference between reversibility of the
complete configuration, equality of entropies and change in the areal reading.
These properties belong to the specified sectoral family;
the physical transport of the bipartition is addressed in
section~\ref{v:v:ant:sec:biparticion-transporte}.

\subsection{First modular law and finite area variations}
\label{v:v:ant:sec:ii-area-ley-modular}

The derivatives of \eqref{v:v:ant:eq:ii-respuesta-entropica} vary the parameter
\(\Delta\) of a family of states. The modular first law, by contrast, fixes
a reference state and considers variations of the observed state
around it. Take
\(\rho_0=\rho_A(\Delta_0)\), with \(\Delta_0\) real and finite,
and keep \(\Delta_0,\gamma_{\HMT},a_0,\ell_P\) fixed.
The branch \(\gamma_{\HMT}>0\), \(a_0>0\) is retained.
The normalized sectoral areas are
\begin{equation}
 \mathcal A_m^{\rm n}=\gamma_{\HMT}a_0N_m,
 \qquad
 \widehat{\mathcal A}_m=\ell_P^2\mathcal A_m^{\rm n},
 \qquad m\in\{90,120\}.
 \label{v:v:ant:eq:ii-area-sector-normalizado}
\end{equation}
The logarithm of the reference state takes the form
\begin{equation}
 K_0=-\log\rho_0
 =c_0I+\beta_{90}\mathcal A_{90}^{\rm n}
         +\beta_{120}\mathcal A_{120}^{\rm n},
 \qquad
 \beta_m=\frac{m\Delta_0}{\gamma_{\HMT}a_0},
 \quad c_0=\log Z(\Delta_0).
 \label{v:v:ant:eq:ii-area-beta-sectorial}
\end{equation}
In particular, \(\beta_{120}=\tfrac43\beta_{90}\); for
\(\Delta_0\ne0\), their quotient is \(4/3\). The coefficients
arise from expressing the same modular operator in areal units.

\begin{theorem}[First modular law in areal coordinates]
\label{v:v:ant:thm:ii-area-primera-ley}
Let \(\sigma(\varepsilon)\) be a differentiable family of normalized
states on an interval containing zero, with
\(\sigma(0)=\rho_0\). Then
\begin{equation}
 \left.\frac{\dd s(\sigma(\varepsilon))}{\dd\varepsilon}\right|_0
 =\sum_{m=90,120}\beta_m
   \left.\frac{\dd\langle\mathcal A_m^{\rm n}\rangle_
                       {\sigma(\varepsilon)}}{\dd\varepsilon}\right|_0.
 \label{v:v:ant:eq:ii-area-primera-ley}
\end{equation}
\end{theorem}
\begin{proof}
The state \(\rho_0\) has full rank. In its neighborhood, the derivative
of the trace of \(-\sigma\log\sigma\) is
\(-\Tr[\dot\sigma(\log\rho_0+I)]\); this follows by diagonalizing
\(\rho_0\) in the formula for the spectral derivative of the trace.
Normalization implies \(\Tr\dot\sigma=0\), so the
derivative is \(\Tr(\dot\sigma K_0)\). Substituting
\eqref{v:v:ant:eq:ii-area-beta-sectorial} also cancels \(c_0I\) and
produces the identity.
\end{proof}

\begin{theorem}[Finite area identity and relative-entropy deficit]
\label{v:v:ant:thm:ii-area-ley-finita}
For any normalized state \(\sigma\) on the same finite space,
define \(\Delta_\sigma\langle B\rangle
=\Tr[(\sigma-\rho_0)B]\). The identity
\begin{equation}
 \boxed{
 s(\sigma)-s(\rho_0)
 =\sum_{m=90,120}\beta_m
      \Delta_\sigma\langle\mathcal A_m^{\rm n}\rangle
   -D(\sigma\Vert\rho_0).}
 \label{v:v:ant:eq:ii-area-ley-finita}
\end{equation}
holds.
Here \(D(\sigma\Vert\rho_0)=\Tr[\sigma(\log\sigma-\log\rho_0)]\)
is finite, nonnegative, and interpreted with \(0\log0=0\).
\end{theorem}
\begin{proof}
By definition,
\(D(\sigma\Vert\rho_0)=-s(\sigma)+\Tr(\sigma K_0)\), whereas
\(s(\rho_0)=\Tr(\rho_0K_0)\). Subtraction gives
\(s(\sigma)-s(\rho_0)=\Delta_\sigma\langle K_0\rangle
-D(\sigma\Vert\rho_0)\). Substituting
\eqref{v:v:ant:eq:ii-area-beta-sectorial} proves the formula because the
traces of both states are one.

To see positivity, diagonalize \(\sigma\) with eigenvalues
\(p_i\) and vectors \(u_i\), and \(\rho_0\) with eigenvalues
\(q_j>0\) and vectors \(v_j\). Put
\(r_i=\langle u_i,\rho_0u_i\rangle\). Concavity of the logarithm
gives
\(\sum_j|\langle u_i,v_j\rangle|^2\log q_j\le\log r_i\).
Hence
\(D(\sigma\Vert\rho_0)\ge\sum_i p_i\log(p_i/r_i)\ge0\);
the last inequality is convexity of \(t\log t\) applied
with weights \(r_i\), which sum to one. Strict positivity of the
\(q_j\) ensures that all reference logarithms are finite.
\end{proof}

The finite identity preserves the informational distance from the
reference state; its tangent term is
\eqref{v:v:ant:eq:ii-area-primera-ley}. For the physical areas, the
corresponding coefficients are \(\beta_m/\ell_P^2\).
Multiplying the equation by the common transducer \(k_B\) publishes
the dimensional entropy. Thus, the relation between geometry and information
remains exact also for finite variations, with its explicit
correction and without modifying the parameters of the reference state.
